% Copyright (c) 2026 Eugene Joseph M. Ragasa.
% SPDX-License-Identifier: Apache-2.0
% Author: Eugene Joseph M. Ragasa
% Formal proof completed: 2026-10-27
\documentclass[11pt]{article}
\usepackage[T1]{fontenc}
\usepackage{amsmath,amssymb}
\usepackage[margin=1in]{geometry}
\usepackage[hidelinks]{hyperref}
\title{Proof Dossier: Singleton-Family Feasibility}
\author{Eugene Joseph M. Ragasa}
\date{Formal proof completed: 2026-10-27}
\begin{document}
\raggedright
\maketitle

\section{Identity and status}
Statement ID \texttt{BG-ALIGNMENT-THM-016}; status \texttt{proved; no admissions}.  Lean declaration
\path{ProjectKoios.BandGauge.Alignment.alignmentFeasible_singleton_iff} in
\path{ProjectKoios/BandGauge/Alignment/Family.lean}.

\section{Exact mathematical statement}
Feasibility in the singleton family $\{C\}$ is equivalent to acceptance of its sole alignment $C$.

\section{Assumptions and conventions}
The alignment and acceptance predicate are arbitrary; no topology or decidability assumption is needed.

\section{Mathematical argument}
In the forward direction, singleton membership identifies the feasible witness with $C$, after which its acceptance proof rewrites to acceptance of $C$.  In the reverse direction, use $C$ itself, its singleton membership, and the supplied acceptance proof.

\section{Lean realization}
The Lean proof follows the mathematical argument directly using the definitions and the
previously proved helper declarations in the same module family.  The declaration was checked
with Lean \texttt{4.35.0} and Mathlib commit \texttt{56a005855ea17e0b0958f73ec801ae140b0a1679}.

\section{Provenance and contribution boundary}
This is source-neutral finite mathematics at Formal merge
\texttt{d9fe61eba430fd2c101dc53d85a2e0978254a718}.  No external manuscript controls the Lean statement, and no priority or
novelty claim is made for the underlying algebra or set theory.

\section{Audit result}
The named axiom report lists \texttt{propext}, \texttt{Classical.choice}, and
\texttt{Quot.sound} and no repository-owned axiom.  The repository admission audit found no
\texttt{sorry} or \texttt{admit} in the checked source.

\section{Boundary and nonclaims}
This is logical normalization of an existential over a singleton, not an optimizer or uniqueness theorem for a larger family.

\section{Review status}
AI-assisted exposition, checked against the Lean source and theorem manifest.  Kernel checking
establishes the exact proposition only.  No independent expert mathematical, scientific,
novelty, or publication review is recorded by this dossier.

\end{document}
